\documentclass[conference]{IEEEtran}
\usepackage[T1]{fontenc}
\usepackage[utf8]{inputenc}
\usepackage{amsmath,amssymb}
\usepackage{array}
\usepackage{booktabs}
\usepackage{multirow}
\usepackage{url}
\usepackage{microtype}
\usepackage{balance}
\title{Can We More Effectively Diagnose and Treat Complex Mental Disorders?\\Mechanism-Guided Adaptive Psychiatry}
\author{Research Plan Author Agent\\Xcientist four-agent simulation}
\begin{document}
\maketitle

\begin{abstract}
Depression, schizophrenia, anxiety, and related presentations overlap in symptoms, biology, treatment response, and course. This report proposes Mechanism-Guided Adaptive Psychiatry (MGAP), a clinician-supervised framework that treats diagnosis and treatment as a repeated inference-and-control problem. MGAP combines transdiagnostic symptom dimensions, multimodal biological features, patient context, treatment history, and time-varying response signals. Its main test is not whether a biomarker correlates with symptoms, but whether a calibrated policy improves treatment matching, durability, and time to an effective clinical review under held-out validation. The design compares diagnosis-only, symptom-dimensional, mechanism, and repeated-update policies. It models inflammation-associated dysconnectivity as a probabilistic subgroup mechanism, uses imaging-genomic features only through governed transformations, and separates association, prediction, treatment moderation, and clinical utility. A treatment-specific longitudinal model and a constrained utility function define the primary analysis. No clinical experiment or observed outcome is claimed; all results are prospective and the protocol is design-only. The central prediction is that repeated, uncertainty-aware, transdiagnostic modeling will improve decisions for a subset of patients without producing a universal psychiatric biomarker. Clinical supervision, subgroup calibration, missing-modality analysis, and a model pause rule are mandatory.
\end{abstract}

\begin{IEEEkeywords}
computational psychiatry, transdiagnostic biomarkers, precision medicine, treatment response, clinical decision support, dysconnectivity, adaptive care
\end{IEEEkeywords}

\section{Introduction}
Complex mental disorders rarely respect the boundaries implied by diagnostic labels. Depression, anxiety, schizophrenia-spectrum conditions, and comorbid presentations can share symptoms while differing in biology, course, context, and treatment response. The practical consequence is familiar: an intervention can help several symptoms at once, fail for another patient, or lose benefit over time. A more effective system must therefore estimate what is active in a particular patient now, not only assign a category at intake.

This report asks whether diagnosis and treatment can be improved by combining transdiagnostic symptom dimensions with multimodal biological features and repeated response measurements. The proposed answer is Mechanism-Guided Adaptive Psychiatry (MGAP). It is a research program, not an autonomous clinical product. A clinician remains responsible for diagnosis, risk management, and treatment authorization. The computational system estimates response trajectories, reports uncertainty, and proposes a review action inside a pre-specified governance policy.

The proposal makes a strong but bounded claim. Psychiatric care should become more effective when it is modeled as repeated inference and control: characterize dimensions and context, estimate plausible mechanisms, select among clinically appropriate paths, monitor early response and tolerability, and adapt when the evidence changes. This is stronger than merely collecting more biomarkers. It requires the model to demonstrate prospective utility, transportability, and safety.

A diagnosis-only model is a necessary negative control. A symptom-dimensional model tests whether better structured phenotyping accounts for most of the gain. A mechanism model asks whether network, inflammatory, physiological, or imaging-genomic features add predictive information. MGAP adds time-varying updates and a clinician-supervised policy. These nested comparisons reveal whether additional biological complexity earns its cost.

The report is organized as a four-agent handoff. The Survey Agent compiles evidence and research gaps. The Idea Agent selects MGAP from competing routes. The ExperimentDesign Agent specifies the cohort, model, endpoints, and falsification rules. The Author Agent freezes those inputs into an IEEE-style proposal. There are no observed clinical results, invented subjects, or individual treatment recommendations.

\section{Evidence-Constrained Survey}
\subsection{Heterogeneity requires transdiagnostic measurement}
McQuaid's review argues that depression has heterogeneous symptom profiles and high comorbidity, making reliable personalized biomarkers difficult to identify~\cite{mcquaid2021}. It motivates biomarkers that map onto shared symptoms or constructs across disorders and emphasizes brain activity, connectivity, neuroendocrine signals, inflammation, and individual context. This directly supports a transdiagnostic feature layer. It does not support the use of one marker as a universal diagnostic test.

Kelly and colleagues describe Research Domain Criteria as an attempt to organize psychiatric research around functional dimensions grounded in neurobiology and observable behavior~\cite{kelly2018}. This provides a useful scaffold for dimensions such as threat, reward, cognitive control, salience, arousal, and interoception. It is a research organization principle, not a reason to discard clinical diagnoses. MGAP retains diagnoses as contextual variables while making dimensional outcomes primary for prediction.

\subsection{Inflammation and network mechanisms}
Goldsmith and colleagues review evidence that inflammatory stimuli can disrupt circuits involved in motivation, motor activity, threat detection, anxiety, interoception, and emotional processing~\cite{goldsmith2023}. Inflammatory biomarkers and acute phase reactants are elevated in a subset of patients with unipolar or bipolar depression, anxiety-related disorders, and schizophrenia, and have been associated with differential treatment response and poor outcomes. The review also emphasizes reproducible measurement of inflammation-associated dysconnectivity.

This evidence supports a mechanism module linking inflammation, connectivity, and symptom dimensions. The word subset is essential. MGAP treats inflammation as a probabilistic feature that may be informative for some patients and irrelevant or misleading for others. The model cannot convert an inflammatory signal into a disease label or infer that an anti-inflammatory intervention will be effective without a suitable treatment test.

\subsection{Imaging genomics and translation}
Chen, Liu, and Calhoun review imaging genomics as a bridge between genetic variation and neurobiological traits with potential relevance to diagnosis and precision medicine~\cite{chen2019}. This supports a governed multimodal layer that connects genomic summaries to imaging-derived phenotypes. It also creates a sharp evidentiary boundary. An association between a genetic feature and a network measure is not a treatment predictor; a treatment predictor is not clinical utility; and utility in one site is not transportability.

MGAP therefore uses imaging-genomic features only after quality control, feature freezing, and nested validation. High-dimensional raw data remain in a controlled environment. The clinical-facing output is a calibrated probability and an explanation at the feature-group level, not a genetic verdict.

\subsection{Treatment response is modality-specific}
Thomas and colleagues studied a transdiagnostic internalizing-psychopathology cohort, compared patients with controls, randomized patients to an SSRI or cognitive behavioral therapy, and examined graph measures as biomarkers and correlates of response~\cite{thomas2020}. Their findings motivate treatment-by-feature interactions and longitudinal network measurement. The study does not establish a universal SSRI or CBT predictor. MGAP uses it as a design precedent for randomized treatment comparison, not as proof that a particular connectome metric should guide care.

\subsection{Adaptive intervention and titration}
Lo and Widge review closed-loop neuromodulation as a possible route to patient-specific psychiatric treatment and identify meaningful titration biomarkers as a central challenge~\cite{lo2017}. Their discussion supports repeated measurement and adaptive control but also argues against unsupervised deployment. MGAP transfers the computational principle of feedback while limiting actions to clinician-authorized review, continuation, intensification, or treatment change.

\begin{table*}[t]
\centering
\caption{Evidence anchors and the claims that MGAP is permitted to make.}
\label{tab:evidence}
\small
\begin{tabular}{p{2.5cm}p{4.0cm}p{4.5cm}p{3.5cm}}
\toprule
Evidence anchor & Directly supported proposition & MGAP design use & Forbidden inference \\
\midrule
McQuaid 2021~\cite{mcquaid2021} & Heterogeneous symptoms, comorbidity, and context complicate biomarker discovery & Transdiagnostic dimensions and context-aware prediction & One biomarker diagnoses all patients \\
Goldsmith et al. 2023~\cite{goldsmith2023} & Inflammation can affect symptom-relevant circuits and response in a subset & Probabilistic inflammation-dysconnectivity mechanism stratum & Inflammation is a universal cause or treatment target \\
Chen et al. 2019~\cite{chen2019} & Imaging genomics can link genetic and neurobiological variation & Governed multimodal feature layer & Association equals utility or causation \\
Thomas et al. 2020~\cite{thomas2020} & Transdiagnostic randomized treatment comparison and modality-specific correlates are feasible & Treatment-by-feature and trajectory tests & Universal SSRI or CBT predictor \\
Lo and Widge 2017~\cite{lo2017} & Closed-loop treatment may improve specificity but titration biomarkers are difficult & Repeated updates with safety supervisor & Autonomous treatment change \\
Kelly et al. 2018~\cite{kelly2018} & Dimensional neurobiology can complement categories & RDoC-inspired phenotype organization & Diagnoses are obsolete \\
\bottomrule
\end{tabular}
\end{table*}

\subsection{Research gap ledger}
Five gaps are accepted. First, it remains unresolved whether multimodal transdiagnostic profiles improve held-out response prediction over diagnosis-only and symptom-only baselines. Second, inflammation and network dysconnectivity may identify subgroups, but subgroup boundaries and transportability are uncertain. Third, response is time-varying and modality-specific, so a baseline-only predictor may miss a clinically useful switching window. Fourth, biomarker association is often reported without a separate test of utility, fairness, and safety. Fifth, the boundary between algorithmic recommendation and clinician-governed adaptation requires an operational protocol.

The survey rejects three stronger claims. It rejects the idea that a single inflammatory, imaging, or genomic measure can diagnose every complex mental disorder. It rejects the idea that an algorithm can replace a psychiatrist or psychologist. It rejects autonomous treatment control in the absence of validated titration biomarkers and clinical governance. These rejections narrow the problem without making the program timid: the proposed trial explicitly tests whether a richer and adaptive policy improves real decisions.

\section{Idea Synthesis}
\subsection{Primary direction: MGAP}
MGAP has four linked layers. The phenotype layer represents symptom dimensions, function, sleep, arousal, treatment history, and context. The mechanism layer represents network connectivity, inflammation, physiology, and imaging-genomic summaries as uncertain features. The response layer predicts treatment-specific trajectories and durability. The policy layer converts predictions into a clinician-supervised review suggestion.

The main scientific hypothesis is:
\begin{quote}
A multimodal, transdiagnostic model that updates from early response and tolerability will improve calibrated treatment matching and time to an effective clinical review over diagnosis-only and symptom-only policies, while its gains will be concentrated in mechanism-defined subgroups rather than universal across patients.
\end{quote}

This hypothesis is deliberately more consequential than a biomarker catalog. A biomarker earns a role only if it improves held-out prediction, survives subgroup and missing-modality checks, and changes decisions without increasing burden or inequity.

\subsection{Competing routes}
A diagnosis-only route uses diagnosis, broad severity, demographics, prior treatment, and site. It represents routine category-level information and is the negative control. A symptom-dimensional route adds structured dimensions and history without biological data. It tests whether measurement of phenotype is sufficient.

A single-biomarker route uses one inflammatory or network feature. It is a high-risk negative control because heterogeneity predicts that one feature will fail to generalize. A fully autonomous policy is rejected for safety. It would blur the distinction between prediction and treatment authorization and would be difficult to audit when the model is uncertain.

MGAP retains the strongest parts of each route. It uses diagnosis as context, symptom dimensions as the clinical interface, multimodal biology as probabilistic mechanism evidence, and repeated monitoring as the feedback channel. It does not promise that all layers are needed for all patients.

\subsection{Predictions}
MGAP predicts improved calibration and decision-curve utility over diagnosis-only prediction. It predicts that the symptom-dimensional model will capture a substantial portion of the gain, making it an important and less costly comparator. It predicts that inflammation-dysconnectivity and network features will moderate response for a subset, not all, patients. It predicts that early repeated response signals will improve timing of clinician review. It predicts that missing-modality handling will reduce performance but preserve utility above the diagnosis-only baseline if uncertainty is propagated.

The model also predicts an important null possibility: biological features may improve retrospective discrimination but fail to improve prospective utility. Under this outcome, the biomarker layer does not earn a clinical role. The distinction protects the program from treating a statistically elegant model as an effective intervention.

\subsection{Why the direction is worth testing}
The same framework can answer whether better care depends primarily on better phenotype measurement, mechanistic stratification, or adaptive timing. A symptom-only win would show that structured assessment is sufficient. A mechanism win would support targeted data collection. An adaptive win would show that the timing of inference matters. A null result would block premature precision-psychiatry deployment and identify which layer failed.

\section{ExperimentDesign}
\subsection{Cohort and data layers}
The target cohort includes adults receiving care for depressive, anxiety-related, schizophrenia-spectrum, or mixed and comorbid presentations. Enrollment records the clinician's working diagnosis without treating it as the only phenotype. Participants are stratified by site and treatment pathway. The unit of inference is the patient nested within site, treatment pathway, and time.

The clinical layer measures anhedonia and motivation, threat and anxiety, cognitive control, psychosis-relevant salience, sleep and arousal, interoceptive distress, functioning, medication history, and previous response. The network layer contains quality-controlled resting-state and task-related connectivity summaries, including fronto-limbic and uncinate-fasciculus measures where acquisition is feasible. The molecular layer contains repeated inflammatory, endocrine, or autonomic measures with assay batch and time-of-day metadata. An imaging-genomic layer connects governed genetic or epigenetic summaries to imaging traits. The response layer records symptoms, functioning, tolerability, adherence, sleep, and clinician-rated risk at each monitoring window.

Every feature receives an availability timestamp. The training pipeline may use a feature only if it would have been available at the decision being simulated. This rule is central: a model that sees future response data is not an adaptive treatment model; it is a leakage artifact.

\subsection{Treatment paths and nested comparators}
The design compares four policies. The diagnosis-only policy uses diagnosis, broad severity, demographics, prior treatment, and site. The symptom-dimensional policy adds transdiagnostic symptom dimensions and history. The mechanism policy adds network, inflammatory, physiological, and imaging-genomic features. MGAP adds repeated response updates, uncertainty, and a clinician-supervised action rule.

Treatment paths may include pharmacotherapy, structured psychotherapy, and clinically approved neuromodulation pathways. The design does not assume one path is superior. Where equipoise and ethics permit, a randomized SSRI-versus-CBT component supplies a treatment contrast for estimating moderation. This follows the logic of a transdiagnostic randomized cohort while avoiding the claim that either modality has a universal predictor~\cite{thomas2020}. Neuromodulation is represented as a separately governed path; the algorithm cannot autonomously stimulate a patient.

\begin{table*}[t]
\centering
\caption{Nested policies, required information, and decision criteria.}
\label{tab:policies}
\small
\begin{tabular}{p{3.0cm}p{4.3cm}p{4.3cm}p{3.0cm}}
\toprule
Policy & Information available at decision time & Main scientific question & Required success condition \\
\midrule
Diagnosis-only & Diagnosis, broad severity, demographics, history, site & How far can category-level care prediction go? & Baseline for every comparison \\
Symptom-dimensional & Diagnosis plus structured dimensions and functioning & Does better phenotype measurement explain improvement? & Calibration gain without biological data \\
Mechanism & Symptom features plus network, inflammatory, physiological, and imaging-genomic features & Do biological mechanisms add treatment-specific information? & Replicated incremental value \\
MGAP & Mechanism features plus repeated response and uncertainty-aware policy & Does adaptive timing improve utility and durability? & Utility gain with safety and equity preserved \\
\bottomrule
\end{tabular}
\end{table*}

\subsection{Response trajectory model}
Let $Y_{i,t,d}$ denote the standardized outcome for patient $i$, time $t$, and symptom or function dimension $d$. A treatment-specific hierarchical trajectory is:
\begin{equation}
Y_{i,t,d}=\alpha_d+u_{i,d}+\beta_{a,d}+\gamma_d^T X_i+\rho_{a,d}^T M_i+\tau_d t+\omega_{a,d}t+\epsilon_{i,t,d}.
\label{eq:trajectory}
\end{equation}
Here $a$ is the treatment pathway, $X_i$ is the symptom and history vector, $M_i$ is the mechanism vector, $u_{i,d}$ is an individual random effect, and $\rho_{a,d}$ estimates treatment-by-mechanism moderation. The interaction $\omega_{a,d}$ permits durability to differ by treatment. The model uses only features available before the predicted time point.

For a binary clinically useful response $R_{i,t}$, the prediction model is:
\begin{equation}
\begin{aligned}
\operatorname{logit}P(R_{i,t}=1)={}&\theta_0+\theta_a+\theta_X^T X_i+\theta_M^T M_i\\
&+\theta_T^T H_{i,t}+\theta_{aM}^T(a\otimes M_i).
\end{aligned}
\label{eq:response}
\end{equation}
The vector $H_{i,t}$ contains early response, tolerability, adherence, and monitoring history. The interaction tests whether different mechanism profiles predict different treatment paths. It is not treated as causal unless treatment assignment is randomized or a valid causal design establishes the needed assumptions.

A response trajectory is more informative than one end-point score. A patient can show early improvement followed by relapse, stable partial improvement, or delayed benefit. MGAP reports slope, asymptote, uncertainty, and durability rather than reducing all of these patterns to a binary responder label.

\subsection{Mechanism strata}
Mechanism strata are estimated with a probabilistic latent class or continuous mixture model. The number of strata is selected by pre-registered predictive criteria and stability across bootstrap samples, not by choosing a clinically attractive number. A patient may have posterior probability distributed across strata. The policy receives that uncertainty instead of forcing a hard biological label.

The principal mechanistic candidate is inflammation-associated dysconnectivity. The evidence indicates that inflammatory signals can affect circuits involved in motivation, threat, anxiety, interoception, and emotion, and that elevated inflammatory markers occur in a subset of people across several psychiatric diagnoses~\cite{goldsmith2023}. The design therefore asks whether inflammation and network features jointly moderate response. It does not treat inflammation as a universal cause or assume that reducing inflammation is beneficial.

\subsection{Missing modalities and measurement quality}
Missing scans, blood assays, or genomic data are expected. The primary analysis marginalizes over missing features using a pattern-aware model and records the reason for missingness. Complete-case analysis is secondary. Performance is reported for complete, partial, and clinically typical data patterns. If missingness is associated with access, severity, or site, the pattern is modeled rather than hidden.

Quality control includes assay batch, scanner, motion, preprocessing version, and time-of-day. Imaging-genomic features are derived inside the training partition and then frozen before held-out evaluation. This prevents a high-dimensional feature layer from silently learning the site or treatment pathway.

\subsection{Clinician-supervised adaptive policy}
At monitoring time $t$, the action $a$ is selected from continue, intensify monitoring, request a clinical review, switch pathway, or augment within the approved care plan. The policy score is:
\begin{equation}
\begin{aligned}
U(a|x_t)={}&p_{\mathrm{resp}}(a|x_t)B_{\mathrm{resp}}
-p_{\mathrm{harm}}(a|x_t)C_{\mathrm{harm}}\\
&-B_{\mathrm{burden}}(a)-\lambda\operatorname{Uncertainty}(a|x_t).
\end{aligned}
\label{eq:utility}
\end{equation}
The benefit, harm cost, burden, and uncertainty penalty are fixed before confirmatory analysis with clinical and patient-partner input. The uncertainty penalty prevents a weakly supported action from outranking a safer review. The output lists dominant feature groups, evidence timestamp, confidence interval, and missingness. A clinician can override the suggestion, and the override reason is analyzed as data.

Closed-loop work motivates the feedback principle but also identifies biomarker titration as difficult~\cite{lo2017}. MGAP therefore limits the policy to decision support. It cannot prescribe, discontinue treatment, or initiate neuromodulation.

\subsection{Causal and predictive separation}
Four claims are separated. Association asks whether a feature relates to an outcome in a specified cohort. Prediction asks whether a model generalizes to held-out patients or sites. Treatment moderation asks whether a feature changes relative response across randomized pathways. Clinical utility asks whether using the model changes care favorably without unacceptable harm.

Inflammation-associated dysconnectivity may be a mechanistic hypothesis, but a cross-sectional association does not prove that inflammation causes a symptom or that an anti-inflammatory treatment will work. The randomized component estimates moderation for its assigned paths. Observational policy evaluation uses explicit target-trial assumptions and reports residual confounding.

\subsection{Validation sequence}
The primary split is time-ordered and site-held-out. Feature selection and tuning occur inside the training partition. External validation uses an independent cohort where possible. The sequence is:
\begin{enumerate}
\item Lock feature timestamps, treatment windows, and leakage exclusions.
\item Fit diagnosis-only and symptom-dimensional baselines.
\item Estimate mechanism strata and treatment interactions.
\item Fit response trajectories and the uncertainty-aware adaptive policy.
\item Validate by site, time period, subgroup, and missingness pattern.
\item Evaluate calibration, decision utility, safety, burden, and clinician overrides.
\item Run simulation-based recovery and sensitivity analyses.
\end{enumerate}

\section{Expected Interpretations and Falsification}
\subsection{Interpretation matrix}
The experiment is designed to produce a decision-relevant result rather than a biomarker leaderboard. Four outcomes are pre-specified.

A clinical utility outcome occurs if MGAP improves calibrated decisions or time to an effective review over the symptom-dimensional policy, with no unacceptable increase in harm, burden, or subgroup inequity. This would support clinician-supervised adaptive psychiatry.

A phenotype-dominant outcome occurs if the symptom-dimensional policy matches MGAP on held-out utility while biological features add little. This would show that structured assessment and repeated clinical data are more valuable than costly biomarker collection for the tested population.

A mechanism-specific outcome occurs if biological features improve treatment moderation and utility for a stable subset, but not universally. This would support targeted data collection and subgroup-specific care rather than one psychiatric biomarker.

A prediction-without-utility outcome occurs if MGAP improves AUC or calibration but does not improve decisions. The model would not earn a clinical role. A model is useful only when the prediction changes care safely and measurably.

\begin{table}[t]
\centering
\caption{Pre-registered interpretation rules for MGAP.}
\label{tab:decisions}
\footnotesize
\begin{tabular}{p{2.0cm}p{2.5cm}p{2.5cm}}
\toprule
Outcome & Statistical signature & Interpretation \\
\midrule
Utility & Held-out utility gain with safety preserved & Adaptive mechanism-guided care is supported \\
Phenotype & Symptom policy matches MGAP & Better measurement dominates biology \\
Mechanism & Subgroup interaction replicates & Targeted biomarker collection is justified \\
No translation & Prediction gain without utility & Do not deploy the added complexity \\
\bottomrule
\end{tabular}
\end{table}

\subsection{Explicit falsification}
MGAP is falsified if nested held-out evaluation shows no clinically meaningful utility improvement over the symptom-dimensional baseline, if mechanism strata are unstable across sites or time, or if early updating fails to predict response better than baseline-only information. It is also falsified as a care policy if recommendations worsen adverse-effect or treatment burden, increase false reassurance, or fail subgroup calibration beyond pre-registered limits.

The inflammation-dysconnectivity module is falsified if its treatment interaction does not replicate after batch, medication, and leakage controls. The imaging-genomic layer is not retained if it adds no external predictive value after feature freezing. The adaptive layer is not retained if clinician review timing is unchanged or if the policy creates unnecessary treatment switches. A null result is a useful decision: it prevents an attractive biomarker story from becoming an ineffective clinical system.

\subsection{Fairness and uncertainty}
Performance is reported by subgroup and missingness pattern. A single average AUC can hide a model that is well calibrated for one group and systematically overconfident for another. The primary fairness report includes calibration intercepts and slopes, false-negative rates for clinically important nonresponse, and decision utility by subgroup. If subgroup estimates are too uncertain, the model reports insufficient evidence rather than producing a precise-looking comparison.

Uncertainty is an input to action. Low confidence should favor a clinical review or additional measurement, not a confident switch. The policy is therefore evaluated on selective prediction: how often it abstains, whether abstention identifies difficult cases, and whether abstention is equitably distributed.

\section{Implementation and Governance}
\subsection{Calibration before clinical use}
The project begins with a retrospective, fully de-identified benchmark and a synthetic recovery study. A prospective observational cohort follows only after the feature schema and leakage rules are frozen. A stepped implementation study is considered only after external validation and independent safety review. This sequence prevents the model from changing care before its errors are characterized.

All preprocessing, feature selection, and tuning occur inside the training partition. A model version is tied to its feature dictionary, assay and scanner versions, and decision threshold. The clinical output is a probability, uncertainty interval, dominant feature group, data timestamp, and suggested review action. It is not a diagnosis and not a prescription.

\subsection{Clinician role and override protocol}
A qualified clinician reviews every model-supported suggestion. The clinician may accept, reject, defer, or request more information. Reasons for override include risk, patient preference, contraindication, missing context, implausible feature, and model uncertainty. Overrides are logged as a safety and usability outcome. The system is not allowed to learn online from a single override without a controlled model update.

High-risk indicators follow established clinical procedures and are not handled by the adaptive model alone. Any medication change, psychotherapy change, or neuromodulation decision requires authorized clinical action. This boundary is a design requirement, not a future aspiration.

\subsection{Data governance}
Mental-health, genomic, and imaging data are sensitive. Access is role-based, clinical identity is separated from model development, and every feature has a provenance record. Participants are told that algorithmic output is investigational. Data sharing uses the minimum necessary resolution and a governed access process. The audit log records model version, input availability, output, clinician decision, and outcome window.

The model pause rule is triggered by predefined safety or equity signals, including an unexpected increase in false negatives, systematic calibration drift, or subgroup harm. When triggered, the system reverts to the clinician's ordinary workflow while the model is reviewed. No deployment claim is made without a monitored rollback plan.

\subsection{Sample size and power}
Sample size is determined by simulation-based design analysis. The simulation varies patient count, site count, treatment imbalance, response prevalence, missing-modality rate, mechanism-stratum separation, treatment interaction, and drift. A design is retained only if it can recover a known treatment-by-mechanism interaction and reject a spurious biomarker under realistic leakage.

The unit of uncertainty is the patient, not a row of a longitudinal table. Repeated observations are nested within patients and sites. A randomized SSRI-versus-CBT component is powered for moderation only after pilot estimates of reliability and response variance. If a treatment pathway is available at one site only, the result is labeled site-limited and excluded from the strongest transportability claim.

\section{Risks and Reproducibility}
\subsection{Construct validity}
Mental disorders contain multiple symptom dimensions and mechanisms. MGAP defines its latent mechanism profiles operationally as reproducible patterns that improve treatment-specific prediction after adjustment for phenotype and context. It does not claim that these profiles are the true biological kinds of mental illness. If a different decomposition proves more stable, raw features and time-stamped outcomes remain usable.

\subsection{Confounding and causal overreach}
Treatment choice in ordinary care is confounded by severity, access, clinician preference, and prior response. Historical policy emulation cannot by itself establish that a biomarker causes response. Randomized treatment contrasts are used for moderation when appropriate. Observational analyses state assumptions and avoid causal language when those assumptions are not credible.

\subsection{Distribution shift}
Scanner models, assay batches, prescribing patterns, health-system access, and symptom composition can drift. Site-held-out and time-held-out evaluation are therefore primary. A model that performs well only at the development site is not a precision-psychiatry success. Monitoring checks drift before using a probability to support a clinical review.

\subsection{Missingness and measurement burden}
Biomarker collection can be expensive or inaccessible. The study reports performance against the number and burden of modalities, not only against predictive accuracy. If a blood assay adds a small gain at large cost while symptom dimensions provide most utility, the simpler policy is preferred. Missingness can encode inequality, so it is modeled as both a technical and social variable.

\subsection{Reproducibility}
The repository contains a frozen feature manifest, data dictionary, preprocessing code, model code, synthetic data, simulation-based calibration, and held-out evaluation scripts. The preregistration specifies outcomes, subgroups, model comparisons, and stopping rules. The project reports negative results, abstentions, clinician overrides, and model pauses. A source manifest binds every literature claim to its evidence capsule.

\subsection{Scope of this report}
This is a research proposal, not a clinical study. It reports no patient outcome, no fitted coefficient, no treatment recommendation, and no claim that a biomarker is clinically validated. Its strong conclusion is methodological and actionable: complex mental disorders should be studied with transdiagnostic dimensions, mechanism-specific hypotheses, and adaptive monitoring, but a biomarker layer must earn its clinical role through held-out utility and safety.

\section{Conclusion}
More effective diagnosis and treatment of complex mental disorders is possible in principle, but the route is not a universal psychiatric biomarker or an autonomous algorithm. The route is a clinician-supervised system that represents the patient at several levels: symptoms and function, plausible biological mechanisms, treatment history, and changing response.

MGAP makes this route testable. It compares diagnosis-only and symptom-dimensional baselines with mechanism and adaptive policies. It models inflammation-associated dysconnectivity as a subset mechanism, imaging-genomic data as a governed translational layer, and treatment response as a trajectory. It asks whether early repeated evidence can improve the timing of a clinical review. It evaluates calibration, utility, burden, subgroup equity, missing modalities, and drift.

The expected result is not that every patient receives a better treatment from a complex model. It is that a subset of heterogeneous patients can receive more timely and better matched care when the system combines dimensions, mechanisms, and feedback. Some patients will be served as well by structured symptom assessment alone. This is not failure; it is a reason to avoid unnecessary data collection.

The decisive standard is clinical utility under uncertainty. If MGAP predicts but does not improve decisions, it should not be deployed. If it improves utility for mechanism-defined subgroups while preserving safety and clinician control, it can become a practical bridge from biological psychiatry to individualized care.

\balance
\begin{thebibliography}{00}
\bibitem{mcquaid2021} R. J. McQuaid, Transdiagnostic biomarker approaches to mental health disorders: Consideration of symptom complexity, comorbidity and context, \emph{Brain, Behavior, and Immunity - Health}, vol. 16, p. 100303, 2021, \url{https://doi.org/10.1016/j.bbih.2021.100303}.
\bibitem{goldsmith2023} D. R. Goldsmith, M. Bekhbat, N. D. Mehta, and J. C. Felger, Inflammation-related functional and structural dysconnectivity as a pathway to psychopathology, \emph{Biological Psychiatry}, vol. 93, no. 5, pp. 405--418, 2023, \url{https://doi.org/10.1016/j.biopsych.2022.11.003}.
\bibitem{chen2019} J. Chen, J. Liu, and V. D. Calhoun, The translational potential of neuroimaging genomic analyses to diagnosis and treatment in the mental disorders, \emph{Proc. IEEE}, vol. 107, no. 5, pp. 912--927, 2019, \url{https://doi.org/10.1109/JPROC.2019.2913145}.
\bibitem{thomas2020} P. J. Thomas et al., Graph theoretical measures of the uncinate fasciculus subnetwork as predictors and correlates of treatment response in a transdiagnostic psychiatric cohort, \emph{Psychiatry Res. Neuroimaging}, vol. 299, p. 111064, 2020, \url{https://doi.org/10.1016/j.pscychresns.2020.111064}.
\bibitem{lo2017} M.-C. Lo and A. S. Widge, Closed-loop neuromodulation systems: next-generation treatments for psychiatric illness, \emph{Int. Rev. Psychiatry}, vol. 29, no. 2, pp. 191--204, 2017, \url{https://doi.org/10.1080/09540261.2017.1282438}.
\bibitem{kelly2018} J. R. Kelly, G. Clarke, J. F. Cryan, and T. G. Dinan, Dimensional thinking in psychiatry in the era of the Research Domain Criteria (RDoC), \emph{Ir. J. Psychol. Med.}, vol. 35, no. 2, pp. 89--94, 2018, \url{https://doi.org/10.1017/ipm.2017.7}.
\end{thebibliography}
\end{document}
